\documentclass{article}
\textheight 23.5cm \textwidth 15.8cm
%\leftskip -1cm
\topmargin -1.5cm \oddsidemargin 0.3cm \evensidemargin -0.3cm
%\documentclass[final]{siamltex}

\usepackage{ctex}
\usepackage{verbatim}
\usepackage{fancyhdr}
\usepackage{graphicx}
\usepackage{booktabs}

%\pagestyle{fancy} \lhead{FDM Homework Template} \chead{}
%\rhead{\bfseries Yan XU}
%
%\lfoot{} \cfoot{} \rfoot{\thepage}
%\renewcommand{\headrulewidth}{0.4pt}
%\renewcommand{\footrulewidth}{0.4pt}
%
\title{第一周程序作业}
\author{孙天阳 SA23001051}
\date{2024年9月25日}
\begin{document}
\maketitle

\section{引言}
本次实验旨在利用有限元方法求解两点边值问题，方程形式如下：
\[
-u'' = f, \quad 0 < x < 1,
\]
其中，边界条件为 \( u(0) = u(1) = 0 \)。通过构造有限元离散方程并采用线性基函数，我们求解了不同网格划分下的数值解，并分析其在 \( L^2 \) 和 \( H^1 \) 范数下的误差收敛性。选择的函数 \( f(x) = 2 \cos(x) - (x-1) \sin(x) \)，其精确解为 \( u(x) = (x - 1) \sin(x) \)。

\section{方法}
为了求解上述边值问题，我们采用了以下步骤：
\begin{itemize}
    \item 首先，将区间 \( [0, 1] \) 均匀划分为 \( N+1 \) 个子区间，生成对应的网格节点和步长 \( h \)。
    \item 其次，组装刚度矩阵 \( A \) 和右端项向量 \( F \)，通过数值积分构造离散方程。
    \item 解方程 \( -A u_h = F \) 获得数值解 \( u_h \)。
    \item 计算数值解与精确解之间的 \( L^2 \) 和 \( H^1 \) 误差，并通过误差的对数关系计算收敛阶数。
\end{itemize}
在具体实现中，针对不同网格密度 \( N = 10, 20, 40, 80 \)，我们分别计算数值解并分析其收敛性。

\section{结果}
\begin{table}[h!]
    \centering
    \caption{\( L^2 \) 和 \( H^1 \) 误差及收敛阶数}
    \label{tab:results}
    \begin{tabular}{ccccc}
        \toprule
        \( N \) & \( L^2 \) 误差 & \( L^2 \) 收敛阶数 & \( H^1 \) 误差 & \( H^1 \) 收敛阶数 \\
        \midrule
        10  & \( 1.40854 \times 10^{-3} \) & -    & \( 4.90167 \times 10^{-2} \) & - \\
        20  & \( 3.86443 \times 10^{-4} \) & 1.87 & \( 2.56657 \times 10^{-2} \) & 0.93 \\
        40  & \( 1.01409 \times 10^{-4} \) & 1.93 & \( 1.31444 \times 10^{-2} \) & 0.97 \\
        80  & \( 2.42843 \times 10^{-5} \) & 2.06 & \( 6.65316 \times 10^{-3} \) & 0.98 \\
        \bottomrule
    \end{tabular}
\end{table}
\section{讨论}
实验中我们计算了不同网格划分下的误差及收敛阶数，结果如表 \ref{tab:results} 所示。可以看到，随着网格密度的增加，\( L^2 \) 和 \( H^1 \) 误差逐渐减小，且收敛阶数总体上接近理论预测的值。具体来说，\( L^2 \) 误差的收敛阶数在 1.87 至 2.06 之间，接近理论上的 2；而 \( H^1 \) 误差的收敛阶数接近 1，这与线性有限元方法的理论结果基本一致。

在编程实现上，我们将整个程序模块化，分解为几个独立的函数，包括生成网格、组装矩阵、求解有限元问题以及计算误差等部分。这样的设计思路不仅使得程序结构更加清晰，易于维护和理解，还为将来进一步拓展和改进提供了灵活性。
\end{document}
